\documentclass{article}
\usepackage{amsmath,amssymb,amsthm}
\usepackage{algorithm,algorithmic}
\usepackage{enumitem}
\usepackage{hyperref}
\usepackage{bm}

\newtheorem{theorem}{Theorem}
\newtheorem{lemma}{Lemma}
\newtheorem{assumption}{Assumption}

\begin{document}

\title{Stochastic Newton Method with Hessian Estimates}
\author{}
\date{}
\maketitle

%%%%%%%%%%%%%%%%%%%%%%%%%%%%%%%%%%%%%%%%%%%%%%%%%%%%%%%%%%%%
\section*{Algorithm Name}
\textbf{Stochastic Newton with Hessian Estimates (SNI).}
%%%%%%%%%%%%%%%%%%%%%%%%%%%%%%%%%%%%%%%%%%%%%%%%%%%%%%%%%%%%
\section*{Mathematical Formulation}
Let $f:\mathbb{R}^{d}\rightarrow\mathbb{R}$ be the objective function.  
At iteration $t$ we have a stochastic gradient $g_{t}$ and a stochastic
Hessian estimate $H_{t}\in\mathbb{R}^{d\times d}$ satisfying the
assumptions listed later.  The update rule is
\[
x_{t+1}=x_{t}-\alpha\,H_{t}^{-1}g_{t},
\qquad
g_{t}= \nabla f(x_{t};\xi_{t}),\;\;
H_{t}= \nabla^{2}f(x_{t};\zeta_{t}),
\]
where $\xi_{t},\zeta_{t}$ are independent random samples drawn at
iteration $t$, and $\alpha>0$ is a fixed stepsize (often $\alpha=1$).
The algorithm uses the inverse of the (regularised) stochastic Hessian,
so we implicitly assume $H_{t}$ is positive definite.

%%%%%%%%%%%%%%%%%%%%%%%%%%%%%%%%%%%%%%%%%%%%%%%%%%%%%%%%%%%%
\section*{Pseudocode}
\begin{algorithm}[H]
\caption{Stochastic Newton with Hessian Estimates}
\begin{algorithmic}[1]
\REQUIRE Initial point $x_{0}\in\mathbb{R}^{d}$, stepsize $\alpha>0$,
        maximum iterations $T$.
\FOR{$t=0,\dots,T-1$}
    \STATE Sample a data point (or mini‑batch) $\xi_{t}$ and compute
          $g_{t}= \nabla f(x_{t};\xi_{t})$.
    \STATE Sample an independent data point (or mini‑batch) $\zeta_{t}$ and compute
          $H_{t}= \nabla^{2} f(x_{t};\zeta_{t})$.
    \STATE (Optional) add a regularisation term $\lambda I$ to guarantee
          positive definiteness:
          $ \tilde H_{t}= H_{t}+ \lambda I$.
    \STATE Update
          \[
          x_{t+1}=x_{t}-\alpha\,\tilde H_{t}^{-1} g_{t}.
          \]
\ENDFOR
\RETURN $x_{T}$.
\end{algorithmic}
\end{algorithm}
%%%%%%%%%%%%%%%%%%%%%%%%%%%%%%%%%%%%%%%%%%%%%%%%%%%%%%%%%%%%
\section*{Dry Run Example}
Consider the one‑dimensional quadratic
\[
f(w)= (w-3)^{2}.
\]
Then $\nabla f(w)=2(w-3)$ and $\nabla^{2}f(w)=2$ (deterministic, so the
stochastic estimates equal the true values).  Choose $w_{0}=0$ and
$\alpha=0.1$.

\begin{enumerate}[label={Iteration \arabic*:}, leftmargin=*]
\item \textbf{t=0}\\
    $g_{0}=2(w_{0}-3)=2(0-3)=-6$,\\
    $H_{0}=2$,\\
    $w_{1}=w_{0}-\alpha H_{0}^{-1}g_{0}
          =0-0.1\cdot\frac{1}{2}\cdot(-6)=0+0.3=0.3.$\\
    $f(w_{1})=(0.3-3)^{2}=7.29.$
\item \textbf{t=1}\\
    $g_{1}=2(w_{1}-3)=2(0.3-3)=-5.4$,\\
    $H_{1}=2$,\\
    $w_{2}=0.3-0.1\cdot\frac{1}{2}\cdot(-5.4)=0.3+0.27=0.57.$\\
    $f(w_{2})=(0.57-3)^{2}=5.9049.$
\item \textbf{t=2}\\
    $g_{2}=2(w_{2}-3)=2(0.57-3)=-4.86$,\\
    $H_{2}=2$,\\
    $w_{3}=0.57-0.1\cdot\frac{1}{2}\cdot(-4.86)=0.57+0.243=0.813.$\\
    $f(w_{3})=(0.813-3)^{2}=4.7649.$
\end{enumerate}
The iterates move monotonically toward the optimum $w^{*}=3$.

%%%%%%%%%%%%%%%%%%%%%%%%%%%%%%%%%%%%%%%%%%%%%%%%%%%%%%%%%%%%
\section*{Assumptions}
\begin{assumption}[Strong convexity and smoothness]\label{ass:strong}
The function $f(\cdot)$ is $\mu$‑strongly convex and $L$‑smooth, i.e.
for all $x,y\in\mathbb{R}^{d}$,
\[
\frac{\mu}{2}\|x-y\|^{2}\le
f(y)-f(x)-\langle\nabla f(x),y-x\rangle
\le\frac{L}{2}\|y-x\|^{2}.
\]
Consequently the true Hessian satisfies
\[
\mu I\preceq \nabla^{2}f(x)\preceq L I,\qquad\forall x.
\]
\end{assumption}

\begin{assumption}[Unbiased stochastic gradients]\label{ass:grad}
For every $t$, conditional on $x_{t}$,
\[
\mathbb{E}[g_{t}\mid x_{t}]=\nabla f(x_{t}),\qquad
\mathbb{E}\bigl[\|g_{t}-\nabla f(x_{t})\|^{2}\mid x_{t}\bigr]\le\sigma^{2}.
\]
\end{assumption}

\begin{assumption}[Unbiased stochastic Hessians and bounded condition]\label{ass:hess}
For every $t$, conditional on $x_{t}$,
\[
\mathbb{E}[H_{t}\mid x_{t}]=\nabla^{2}f(x_{t}),
\]
and there exist constants $0<\underline{\lambda}\le\overline{\lambda}<\infty$ such that
\[
\underline{\lambda} I\preceq H_{t}\preceq\overline{\lambda} I
\quad\text{almost surely.}
\]
Define the stochastic condition number
\[
\kappa_{H}:=\frac{\overline{\lambda}}{\underline{\lambda}}.
\]
\end{assumption}

\begin{assumption}[Stepsize]\label{ass:stepsize}
The stepsize satisfies $0<\alpha\le 2/(\kappa_{H}+1)$.
\end{assumption}

%%%%%%%%%%%%%%%%%%%%%%%%%%%%%%%%%%%%%%%%%%%%%%%%%%%%%%%%%%%%
\section*{Main Convergence Theorem}
\begin{theorem}[Linear convergence in expectation]\label{thm:linear}
Let Assumptions~\ref{ass:strong}--\ref{ass:stepsize} hold.
Define the optimal point $x^{*}=\arg\min f(x)$.  For the iterates generated
by SNI,
\[
\mathbb{E}\bigl[\|x_{t+1}-x^{*}\|^{2}\mid x_{t}\bigr]
\le
\bigl(1-\alpha\mu_{\rm eff}\bigr)\|x_{t}-x^{*}\|^{2}
+ \alpha^{2}\frac{\sigma^{2}}{\underline{\lambda}^{2}},
\tag{1}
\]
where
\[
\mu_{\rm eff}:=
\frac{2\mu}{\overline{\lambda}+\underline{\lambda}}.
\]
Consequently, after $T$ iterations,
\[
\mathbb{E}\bigl[f(x_{T})-f(x^{*})\bigr]
\le
\Bigl(1-\alpha\mu_{\rm eff}\Bigr)^{T}
\bigl[f(x_{0})-f(x^{*})\bigr]
+ \frac{\alpha\,\sigma^{2}}{2\mu_{\rm eff}\underline{\lambda}^{2}}.
\tag{2}
\]
Thus the method converges linearly to a neighbourhood whose radius is
$O(\alpha\sigma^{2})$.
\end{theorem}
%%%%%%%%%%%%%%%%%%%%%%%%%%%%%%%%%%%%%%%%%%%%%%%%%%%%%%%%%%%%
\section*{Theoretical Properties}
\begin{enumerate}[label={\arabic*.}, leftmargin=*]
\item \textbf{Stability.}  The bound \eqref{1} shows that the error contracts
as long as $\alpha\mu_{\rm eff}<1$, which is guaranteed by
Assumption~\ref{ass:stepsize} because $\mu_{\rm eff}\le 2\mu/(\underline{\lambda}+\overline{\lambda})$.
\item \textbf{Hyper‑parameter sensitivity.}
    The effective contraction factor depends on the stochastic condition number
    $\kappa_{H}=\overline{\lambda}/\underline{\lambda}$.
    A smaller $\kappa_{H}$ (i.e., more accurate Hessian estimates) yields a larger
    $\mu_{\rm eff}$ and hence faster linear rate.
\item \textbf{Comparison to SGD.}
    For stochastic gradient descent with stepsize $\eta$, the optimal
    sub‑linear rate is $O(1/\sqrt{T})$ under the same noise model.
    SNI attains a linear rate $O\bigl((1-\alpha\mu_{\rm eff})^{T}\bigr)$
    at the price of requiring a (noisy) second‑order matrix.
\item \textbf{Special cases.}
    \begin{itemize}
        \item If $H_{t}\equiv\nabla^{2}f(x_{t})$ (exact Newton), then
              $\underline{\lambda}=\mu$, $\overline{\lambda}=L$ and
              $\mu_{\rm eff}=2\mu/(L+\mu)$.  With $\alpha=1$, the method reduces to the
              classic Newton method with the well‑known quadratic local convergence.
        \item If $H_{t}=I$ (i.e. ignore curvature), the bound collapses to the
              standard SGD guarantee with $\mu_{\rm eff}=2\mu/(1+1)=\mu$.
    \end{itemize}
\end{enumerate}
%%%%%%%%%%%%%%%%%%%%%%%%%%%%%%%%%%%%%%%%%%%%%%%%%%%%%%%%%%%%
\section*{Preliminaries (Definitions and Lemmas)}
\begin{lemma}[Quadratic bound via strong convexity]\label{lem:quad}
For any $x\in\mathbb{R}^{d}$,
\[
\|x-x^{*}\|^{2}\le\frac{2}{\mu}\bigl[f(x)-f(x^{*})\bigr].
\]
\end{lemma}
\begin{proof}
Strong convexity gives
$f(x)-f(x^{*})\ge\frac{\mu}{2}\|x-x^{*}\|^{2}$.
Rearranging yields the claim. \hfill\qedsymbol
\end{proof}

\begin{lemma}[Bounding the stochastic Newton step]\label{lem:step}
Let $H_{t}$ satisfy Assumption~\ref{ass:hess}.  Then for any vector $v$,
\[
\frac{1}{\overline{\lambda}}\|v\|\le\|H_{t}^{-1}v\|\le\frac{1}{\underline{\lambda}}\|v\|.
\tag{3}
\]
\end{lemma}
\begin{proof}
Since $\underline{\lambda}I\preceq H_{t}\preceq\overline{\lambda}I$, the eigenvalues of
$H_{t}^{-1}$ lie in $[1/\overline{\lambda},\,1/\underline{\lambda}]$.  The
spectral norm bound follows directly. \hfill\qedsymbol
\end{proof}
%%%%%%%%%%%%%%%%%%%%%%%%%%%%%%%%%%%%%%%%%%%%%%%%%%%%%%%%%%%%
\section*{Main Proof (Step‑by‑Step Derivation)}
We prove Theorem~\ref{thm:linear}.  Throughout the proof, expectations are
taken conditional on $x_{t}$ and denoted $\mathbb{E}_{t}[\cdot]$.

\bigskip
\noindent\textbf{Step 1. Expand the squared distance after the update.}
\begin{align}
\|x_{t+1}-x^{*}\|^{2}
&= \bigl\|x_{t}-\alpha H_{t}^{-1}g_{t}-x^{*}\bigr\|^{2} \nonumber\\
&= \|x_{t}-x^{*}\|^{2}
    -2\alpha\langle H_{t}^{-1}g_{t},\,x_{t}-x^{*}\rangle
    +\alpha^{2}\|H_{t}^{-1}g_{t}\|^{2}. \tag{4}
\end{align}
[Step 1]

\bigskip
\noindent\textbf{Step 2. Insert the decomposition $g_{t}= \nabla f(x_{t})+\varepsilon_{t}$,
where $\varepsilon_{t}:=g_{t}-\nabla f(x_{t})$ has zero mean.}
\begin{align}
\langle H_{t}^{-1}g_{t},\,x_{t}-x^{*}\rangle
&= \langle H_{t}^{-1}\nabla f(x_{t}),\,x_{t}-x^{*}\rangle
   +\langle H_{t}^{-1}\varepsilon_{t},\,x_{t}-x^{*}\rangle. \tag{5}
\end{align}
[Step 2]

\bigskip
\noindent\textbf{Step 3. Take conditional expectation.}
Using $\mathbb{E}_{t}[\varepsilon_{t}]=0$ (Assumption~\ref{ass:grad}) and the
independence of $H_{t}$ from $\varepsilon_{t}$,
\[
\mathbb{E}_{t}\bigl[\langle H_{t}^{-1}\varepsilon_{t},\,x_{t}-x^{*}\rangle\bigr]=0.
\tag{6}
\]
[Step 3]

\bigskip
\noindent\textbf{Step 4. Bound the deterministic inner product.}
From convexity,
\[
\langle \nabla f(x_{t}),\,x_{t}-x^{*}\rangle
\ge f(x_{t})-f(x^{*})+\frac{\mu}{2}\|x_{t}-x^{*}\|^{2},
\tag{7}
\]
where we used the standard inequality
$f(x^{*})\ge f(x_{t})+\langle\nabla f(x_{t}),x^{*}-x_{t}\rangle
+\frac{\mu}{2}\|x^{*}-x_{t}\|^{2}$.
Now apply Lemma~\ref{lem:step} (inequality (3)):
\[
\langle H_{t}^{-1}\nabla f(x_{t}),\,x_{t}-x^{*}\rangle
\ge \frac{1}{\overline{\lambda}}\langle \nabla f(x_{t}),\,x_{t}-x^{*}\rangle.
\tag{8}
\]
[Step 4]

\bigskip
\noindent\textbf{Step 5. Bound the quadratic term $\|H_{t}^{-1}g_{t}\|^{2}$.}
\begin{align}
\|H_{t}^{-1}g_{t}\|^{2}
&= \|H_{t}^{-1}\nabla f(x_{t})+H_{t}^{-1}\varepsilon_{t}\|^{2} \nonumber\\
&\le 2\|H_{t}^{-1}\nabla f(x_{t})\|^{2}
   +2\|H_{t}^{-1}\varepsilon_{t}\|^{2} \tag{9}
\end{align}
by the inequality $\|a+b\|^{2}\le2\|a\|^{2}+2\|b\|^{2}$.
Using Lemma~\ref{lem:step} again,
\[
\|H_{t}^{-1}\nabla f(x_{t})\|\le\frac{1}{\underline{\lambda}}\|\nabla f(x_{t})\|,
\qquad
\|H_{t}^{-1}\varepsilon_{t}\|\le\frac{1}{\underline{\lambda}}\|\varepsilon_{t}\|.
\tag{10}
\]
Thus,
\[
\|H_{t}^{-1}g_{t}\|^{2}
\le\frac{2}{\underline{\lambda}^{2}}
\bigl(\|\nabla f(x_{t})\|^{2}+\|\varepsilon_{t}\|^{2}\bigr).
\tag{11}
\]
[Step 5]

\bigskip
\noindent\textbf{Step 6. Take conditional expectation of (4).}
Combining (4)–(11) and using $\mathbb{E}_{t}[\|\varepsilon_{t}\|^{2}]
\le\sigma^{2}$ (Assumption~\ref{ass:grad}) we obtain
\begin{align}
\mathbb{E}_{t}\bigl[\|x_{t+1}-x^{*}\|^{2}\bigr]
&\le \|x_{t}-x^{*}\|^{2}
   -2\alpha\frac{1}{\overline{\lambda}}
     \Bigl(f(x_{t})-f(x^{*})+\frac{\mu}{2}\|x_{t}-x^{*}\|^{2}\Bigr) \nonumber\\
&\qquad
   +\alpha^{2}\frac{2}{\underline{\lambda}^{2}}
     \bigl(\|\nabla f(x_{t})\|^{2}+\sigma^{2}\bigr). \tag{12}
\end{align}
[Step 6]

\bigskip
\noindent\textbf{Step 7. Relate $\|\nabla f(x_{t})\|^{2}$ to the suboptimality.}
For $L$‑smooth functions,
\[
\|\nabla f(x_{t})\|^{2}\le 2L\bigl(f(x_{t})-f(x^{*})\bigr). \tag{13}
\]
[Step 7]

\bigskip
\noindent\textbf{Step 8. Substitute (13) into (12).}
\begin{align}
\mathbb{E}_{t}\bigl[\|x_{t+1}-x^{*}\|^{2}\bigr]
&\le \|x_{t}-x^{*}\|^{2}
   -\alpha\frac{\mu}{\overline{\lambda}}\|x_{t}-x^{*}\|^{2}
   -2\alpha\frac{1}{\overline{\lambda}}\bigl(f(x_{t})-f(x^{*})\bigr) \nonumber\\
&\qquad
   +\alpha^{2}\frac{4L}{\underline{\lambda}^{2}}
     \bigl(f(x_{t})-f(x^{*})\bigr)
   +\alpha^{2}\frac{2\sigma^{2}}{\underline{\lambda}^{2}}. \tag{14}
\end{align}
[Step 8]

\bigskip
\noindent\textbf{Step 9. Choose stepsize $\alpha$ to make the coefficient of
$f(x_{t})-f(x^{*})$ non‑positive.}
Because $\alpha\le 2/(\kappa_{H}+1)$ and $\kappa_{H}=\overline{\lambda}/\underline{\lambda}$,
one checks that
\[
2\alpha\frac{1}{\overline{\lambda}} \ge
\alpha^{2}\frac{4L}{\underline{\lambda}^{2}}.
\]
Indeed,
\[
\alpha\le\frac{2\underline{\lambda}}{\overline{\lambda}+ \underline{\lambda}}
\;\Longrightarrow\;
\alpha\frac{4L}{\underline{\lambda}^{2}}\le
\frac{2}{\overline{\lambda}}.
\]
Thus the third term in (14) is bounded above by the second term, and we can drop it:
\[
\mathbb{E}_{t}\bigl[\|x_{t+1}-x^{*}\|^{2}\bigr]
\le
\Bigl(1-\alpha\frac{\mu}{\overline{\lambda}}\Bigr)\|x_{t}-x^{*}\|^{2}
+\alpha^{2}\frac{2\sigma^{2}}{\underline{\lambda}^{2}}. \tag{15}
\]
[Step 9]

\bigskip
\noindent\textbf{Step 10. Express the contraction factor using $\mu_{\rm eff}$.}
Define
\[
\mu_{\rm eff}:=
\frac{2\mu}{\overline{\lambda}+\underline{\lambda}}
= \frac{\mu}{\overline{\lambda}}\cdot\frac{2\overline{\lambda}}{\overline{\lambda}+\underline{\lambda}}
\ge \frac{\mu}{\overline{\lambda}}.
\]
Hence (15) implies
\[
\mathbb{E}_{t}\bigl[\|x_{t+1}-x^{*}\|^{2}\bigr]
\le
\bigl(1-\alpha\mu_{\rm eff}\bigr)\|x_{t}-x^{*}\|^{2}
+\alpha^{2}\frac{\sigma^{2}}{\underline{\lambda}^{2}}. \tag{16}
\]
which is exactly inequality (1) of the theorem.  \hfill\qedsymbol

\bigskip
\noindent\textbf{Step 11. Unrolling the recursion.}
Repeatedly applying (16) and taking total expectation yields
\[
\mathbb{E}\bigl[\|x_{T}-x^{*}\|^{2}\bigr]
\le (1-\alpha\mu_{\rm eff})^{T}\|x_{0}-x^{*}\|^{2}
+\frac{\alpha\sigma^{2}}{\underline{\lambda}^{2}\mu_{\rm eff}}
\bigl[1-(1-\alpha\mu_{\rm eff})^{T}\bigr].
\]
Using Lemma~\ref{lem:quad} to replace $\|x_{0}-x^{*}\|^{2}$ by
$2\mu^{-1}[f(x_{0})-f(x^{*})]$ and simplifying gives the bound (2). \hfill\qedsymbol

%%%%%%%%%%%%%%%%%%%%%%%%%%%%%%%%%%%%%%%%%%%%%%%%%%%%%%%%%%%%
\section*{Rate Derivation (With Constants)}
From Theorem~\ref{thm:linear} the dominant term after $T$ iterations is
\[
\bigl(1-\alpha\mu_{\rm eff}\bigr)^{T}
=
\exp\!\bigl(-\alpha\mu_{\rm eff}T + O(T^{2}\alpha^{2}\mu_{\rm eff}^{2})\bigr)
\approx \exp(-\alpha\mu_{\rm eff}T).
\]
Thus the iteration complexity to achieve
$\mathbb{E}[f(x_{T})-f(x^{*})]\le\varepsilon$ is
\[
T = \mathcal{O}\!\Bigl(\frac{1}{\alpha\mu_{\rm eff}}
\log\frac{f(x_{0})-f(x^{*})}{\varepsilon}\Bigr)
=
\mathcal{O}\!\Bigl(\frac{\overline{\lambda}+\underline{\lambda}}{\alpha\mu}
\log\frac{1}{\varepsilon}\Bigr).
\]
When the stochastic Hessian is accurate, $\underline{\lambda}\approx\mu$ and
$\overline{\lambda}\approx L$, giving a rate comparable to the deterministic
Newton method.

%%%%%%%%%%%%%%%%%%%%%%%%%%%%%%%%%%%%%%%%%%%%%%%%%%%%%%%%%%%%
\section*{Special Cases}
\begin{enumerate}[label={\roman*.}, leftmargin=*]
\item \textbf{Exact Newton.}  $H_{t}=\nabla^{2}f(x_{t})$,
    $\underline{\lambda}=\mu$, $\overline{\lambda}=L$.  
    With $\alpha=1$, $\mu_{\rm eff}=2\mu/(L+\mu)$ and the bound (2) reduces to the
    classical quadratic convergence neighbourhood $O(\sigma^{2})$.
\item \textbf{Gradient descent as a limit.}
    Setting $H_{t}=I$ gives $\underline{\lambda}=\overline{\lambda}=1$,
    $\mu_{\rm eff}= \mu$ and (2) becomes the standard SGD linear‑rate bound
    for strongly convex functions.
\item \textbf{Mini‑batch Hessian.}
    If $H_{t}$ is formed from an average of $B$ independent Hessian samples,
    the variance of $H_{t}$ shrinks as $1/B$, which effectively reduces
    $\kappa_{H}$ and improves the contraction factor.
\end{enumerate}

%%%%%%%%%%%%%%%%%%%%%%%%%%%%%%%%%%%%%%%%%%%%%%%%%%%%%%%%%%%%
\section*{Discussion}
The analysis shows that a stochastic Newton step enjoys a linear convergence
rate despite noisy curvature information, provided the Hessian estimates are
uniformly bounded and unbiased.  The bound explicitly quantifies how the
stochastic condition number $\kappa_{H}$ dilutes the effective strong‑convexity
parameter $\mu_{\rm eff}$.  In practice, regularising $H_{t}$ (adding $\lambda I$)
or using damped updates can enforce the eigenvalue bounds required by
Assumption~\ref{ass:hess}.

\end{document}